\section{Related Works}

\subsection{Data Synthesis for Simultaneous Speech Translation}

Training SimulST models requires aligned trajectories pairing partial
source speech with partial target translations. Existing approaches
synthesize such data from offline corpora using different strategies to
determine when each target token should be committed.

One line of work constructs prefix-aligned supervision through explicit
alignment. Prefix Alignment \citep{kano-etal-2022-simultaneous} extracts
alignments between bilingual prefix pairs, uses them to segment streaming
input, and fine-tunes a translation model on the resulting prefix-to-prefix
examples. This strategy is attractive because it is simple, interpretable, and
does not require LLM calls, but its deterministic alignment rule cannot
directly represent cases where the safety of a target prefix depends on
unobserved future context.

Simul-MuST-C \citep{Simul-must-c} prompts GPT-4o to re-segment offline
translations into minimal monotonic chunks following source word order,
inspired by the salami technique used by human interpreters.
EAST \citep{EAST} uses GPT-4 to segment transcripts into semantically
coherent chunks under multiple latency settings, reformatted as
interleaved read/write sequences. Both approaches determine commit points
from the current source prefix alone, without verifying whether the
chosen translation remains valid given future context.

InfiniSST \citep{ouyang2025infinisst} constructs trajectories through a
pipeline of forced speech--transcript alignment (MFA) followed by
word-level transcript--translation alignment with SimAlign
\citep{jalili-sabet-etal-2020-simalign}, committing each translation
token once its aligned source chunk arrives. Hibiki takes a softer
approach, using the perplexity of an MT system to identify per-word
optimal delays --- a target word is committed once it becomes predictable
from the current source. While more adaptive than hard alignment, both
methods rely on local signals and do not explicitly reason over what
future source words might require.

TAF and concurrent future-anticipation work
\citep{ouyang2024anticipating} address future dependency by sampling
multiple LLM-predicted source continuations, translating each, and
selecting the longest common literal prefix across hypotheses. These
methods demonstrate that future-aware reasoning substantially improves
the quality--latency tradeoff, but their literal string aggregation is
sensitive to surface variation among hypotheses, and they require LLM
calls at inference time, adding cost and latency to deployment.

In contrast, our method replaces literal prefix matching with token-level distributional consensus during training-data construction, decoupling future-aware reasoning from inference time.
